\documentclass[leqno, a4paper, 12pt]{jsarticle}
\usepackage{amssymb}
\usepackage{amsthm}
\usepackage{amsmath}
\usepackage{comment}
\usepackage{enumitem}
\usepackage{url}

%\usepackage{showkeys}
	%可換図式用
		%\usepackage[all]{xy}
	%重くなるので使わいときはコメント化しておく。
%定理環境 thmはあまり使わずpropをなるべく使う。
%主張は基本的に証明の中で使い、そのため番号を付けない。
%注意環境を付けてみた。
%
\usepackage{}
\newtheorem{thm}{定理}
\newtheorem{prop}[thm]{命題}
\newtheorem{cor}[thm]{系}
\newtheorem{lem}[thm]{補題}
\newtheorem{df}[thm]{定義}
\newtheorem{exam}[thm]{例}
\newtheorem{fact}[thm]{事実}
\newtheorem*{claim}{主張}
\newtheorem*{rmk}{注意}
\renewcommand{\proofname}{{\bf 証明}}
%矢印たち。
%\toは\rightarrowと同じ。\getsは\leftarrowと同じ。同値は\iff
%上向き矢はuparrow逆はdownarrow
\newcommand{\To}{\Longrightarrow}
\newcommand{\Gets}{\Longleftarrow}
%黒板太字
\newcommand{\nn}{\mathbb{N}}
\newcommand{\zz}{\mathbb{Z}}
\newcommand{\qq}{\mathbb{Q}}
\newcommand{\rr}{\mathbb{R}}
\newcommand{\cc}{\mathbb{C}}
%無限大
\newcommand{\mgn}{\infty}
%ギリシャン
\newcommand{\dpst}{\displaystyle}
\newcommand{\ep}{\varepsilon}
\newcommand{\al}{\alpha}
\newcommand{\be}{\beta}
\newcommand{\de}{\delta}
\newcommand{\la}{\lambda}
\newcommand{\La}{\Lambda}
\newcommand{\ga}{\gamma}
\newcommand{\lil}{\la\in \La}
%商集合
\newcommand{\qt}[2]{#1/\! #2}
%enumerate環境
\renewcommand{\labelenumi}{{\rm (\arabic{enumi})}}
%以降alignじゃなくてenumerate を使え。
%なんか演算
\DeclareMathOperator{\card}{card}
\DeclareMathOperator{\supp}{supp}
%なんか演算２
\DeclareMathOperator{\bd}{\partial}
\DeclareMathOperator{\cl}{CL}
\DeclareMathOperator{\nai}{Int}
\DeclareMathOperator{\di}{diam}

\newcommand{\ind}{\mathrm{ind}}

\newcommand{\lind}{\mathrm{Ind}}

\DeclareMathOperator{\emb}{Emb}
%差集合は\setminusで統一する。等号の否定は\neq
%真部分集合は\subsetneqq　偏微分は\partial
%ディスプレイスタイル。\dfracでディスプレイ分数
%空集合は\Oを使う！！！！！！！！！！！！


\DeclareMathOperator{\rank}{rank}


%%%%%%%%%%%%%%%%%

\newcommand{\mydisc}[1]{\mathbf{D}^{#1}}
\newcommand{\mysphere}[1]{\mathbf{S}^{#1}}
\newcommand{\mycube}[1]{\mathbf{I}^{n}}
\newcommand{\mycubesidep}[2]{\mathbf{I}^{#1, #2}_{+}}
\newcommand{\mycubesidem}[2]{\mathbf{I}^{#1, #2}_{-}}

\newcommand{\mycubebd}[1]{\partial\mathbf{I}^{#1}}


\newcommand{\myhemispp}[1]{\mathbf{H}_{+}^{#1}}
\newcommand{\myhemispm}[1]{\mathbf{H}_{-}^{#1}}
\newcommand{\myeqsp}{\mathbf{E}}

\newcommand{\myhl}[2]{H_{#2}(#1)}


%%%%%%%%%%%%%%%%%%


\title{
廉価版領域不変性定理の証明
}
\author{ナンブキトラ}
\date{\today}
			\begin{document}
\maketitle

この文章では領域不変性定理を
次元論などの
難しいことをなるべく使わないで証明してみる．
次元論を捨てたことにより球面への連続写像の拡張で
次元を特徴づける定理の範囲がユークリッド空間に狭まった．
しかし，ユークリッド空間では次元論の代わりに
ユークリッド空間特有の現象，つまり
$\dagger$解析学$\dagger$
を使うことができる(実際にはリプシッツ性があれば今回使う程度のサードの定理は証明できるので，本当に解析っぽい部分はサードの定理で測度ゼロの概念を用いるところのみである)．


以下の証明方法，特にサードの定理を使う部分については
テレンス・タオ\cite{tao}がブログ, 
もしくはその元ネタである
\cite{Kul}
に影響を受けた．

証明の方針は
\cite{dempa3}
と全く同じである．
違うのは
定理
\ref{thm:26thismainthm}
の適応範囲をユークリッド空間の球面のみに絞って
ユークリッド空間特有の性質で証明した部分のみである．
第二節の領域不変性定理の証明は
\cite{dempa3}
をそのまま流用している．

\section{本文}

記号
$U(a, r)$
で$a$を中心とする
半径$r$の開球を表すとする．
\begin{thm}\label{thm:26kinji}
$n, k\in \zz_{\ge 1}$
とし
$K$を
$\rr^{n}$
のコンパクト部分集合とし
$f\colon K\to \rr^{k}$
を
連続写像とする．
このとき
任意の
$\epsilon\in (0, \infty)$
に対して
$C^{\infty}$級
関数
$h\colon \rr^{n}\to \rr^{k}$
が存在して
$\sup_{x\in K}\lVert f(x)-h(x)\rVert<\epsilon$
が成り立つ．
\end{thm}
\begin{proof}
$K$のコンパクト性から
$f$は$K$上一様連続である．
そして
$f$の一様連続性から
ある$\delta\in (0, \infty)$
が存在して
$x, y\in K$
が$\lVert x-y \rVert<\delta$
を満たすならば
$\lVert f(x)-f(y)\rVert<\epsilon$
となる．
そして
$K$のコンパクト性から
ある点列$a_{1}, \dots, a_{m}\in K$
が存在して
$U(a_{1}, \delta), \dots, U(a_{m}, \delta), X\setminus K$
が$\rr^{k}$の開被覆となる．
よって
滑らかな
単位の分割$\{u_{i}\}_{i\in \{1, \dots, m+1\}}$が
存在して
$u_{i}$は$U(a_{i}, \delta)$に従属し
$u_{m+1}$
は$X\setminus K$
に従属する．
よって$x\in K$
のとき
$\sum_{i=1}^{m}u_{i}(x)=1$
となる．
そして
$h\colon \rr^{n}\to \rr^{k}$
を
\[
h(x)=\sum_{i=1}^{m}f(a_{i})\cdot u_{i}(x)
\]
と定義する．
すると$h$は滑らかであり任意の
$x\in K$について
\begin{align*}
\lVert f(x)-h(x) \rVert
&=
\left\lVert \sum_{i=1}^{m}f(x)u(x)-\sum_{i=1}^{m}f(a_{i})u(x)
\right\rVert
=
\left\lVert \sum_{i=1}^{m}(f(x)-f(a_{i}))u(x)
\right\rVert\\
&\le 
\sum_{i=1}^{m}u(x)\left\lVert f(x)-f(a_{i})
\right\rVert
< \sum_{i=1}^{m}u(x)\eta=\eta
\end{align*}
\end{proof}
なので，定理が成り立つことがわかる．
\begin{rmk}
他にも方法はある．まず
関数$f$に対してティーチェの拡張定理を用いて
$f$を$g\colon \rr^{n}\to \rr^{k}$
に拡張する．そしてこの$g$を
例えば多項式近似定理や
ストーン・ワイエルストラスの定理を使って$K$上で近似すると証明できる．
あるいは
軟化子でほぐしてあげると
定理が証明できる．
実は微分可能性はサードの定理を適応する時にしか使わず，
さらに今回使うサードの定理のためには
リプシッツ性があれば十分なので，
上の定理で滑らかな単位の分割の代わりに
リプシッツな単位の分割をこさえてやればこの文書の話は回すことができる．
\end{rmk}


\begin{thm}[Sardの定理スペシフィック]\label{thm:26sard}
関数$f\colon \rr^{n}\to \rr^{n+1}$
は$C^{\infty}$級とする．
このとき
$f(\rr^{n})$は$\rr^{n+1}$
の中で測度ゼロである．
\end{thm}
\begin{proof}
大まかにいえば$f$が局所リプシッツであることから
$f$の像が測度ゼロになることが従う．
つまり
$f$が局所的に$M$リプシッツだとすると
$f$による半径$r$のボールの像は半径$Mr$のボールに含まれる．なのでこのような写像で測度ゼロ性は保たれるのである．
\end{proof}

\begin{lem}\label{lem:retractnameraka}
$n\in \zz_{\ge 0}$で$r\in (0, 1)$とする．
そして
$A_{r}=\{x\in \rr^{n+1}\mid 1-r<\lVert x\rVert<1+r\}$と
おく．
なめらかなレトラクション
$\pi \colon \rr^{n+1}\setminus \{0\}\to \mysphere{n}$
を
\[
\pi(x)=\frac{x}{\lVert x\rVert}
\]
と定義する．
このとき
$r$に依存する定数
$M_{r}$が存在して
$\pi$
は$A_{r}$
上で$M_{r}$リプシッツ写像である．
\end{lem}
\begin{proof}
$A_{r}$は相対コンパクトで
$\pi$の微分のノルムが上から抑えられるからである．
\end{proof}

\begin{cor}\label{cor:extspnameraka}
$n\in \zz_{\ge 0}$
とし
$X$を
$\rr^{n}$
のコンパクト部分集合とし
$A$を$X$の閉部分集合とする．
このとき
任意の$\epsilon \in (0, \infty)$
と
任意の連続写像
$f\colon A\to \mysphere{n}$
に対して
(なめらかな)関数
$h\colon \rr^{n}\to \mysphere{n}$
が存在し
以下の条件を満たす．
\begin{enumerate}
\item 
$\sup_{x\in A}\lVert f(x) -h(x) \rVert<\epsilon$. 
\item 
像$h(\rr^{n})$は$0\in \rr^{n+1}$
を含まない．
\end{enumerate}
\end{cor}
\begin{proof}
$r=\min\{2^{-1}, \epsilon\}$かつ
$\eta=\frac{\min\{\epsilon, r\}}{1000(M_{r}+1)}$
と定義する．
写像
$f\colon A\to \mysphere{n}$ 
を
$f\colon A\to \rr^{n+1}$
とみなす．
そして
定理
\ref{thm:26kinji}
を適応して
なめらかな関数
$u\colon \rr^{n}\to \rr^{n+1}$
で
$\sup_{x\in A}\lVert f(x)-u(x) \rVert<\eta$
となるもの
を得る．
このとき
定理
\ref{thm:26sard}
から
$u(\rr^{n})$
は$\rr^{n+1}$
のなかで
測度ゼロなので
特に$\rr^{n+1}\setminus u(\rr^{n})$
は稠密である．
よって
十分小さい$\alpha\in (0, \epsilon/10)$
を使って$v(x)=u(x)+\alpha$と定義される関数
$v\colon \rr^{n}\to \rr^{n+1}$は
$0\not\in v(\rr^{n})$
を満たす．
このとき
$\sup_{x\in A}\lVert f(x)-v(x)\lVert\le 2\eta$
である．
そして
$h=\pi\circ v\colon \rr^{n}\to \mysphere{n}$
と定義する．
すると
$x\in A$のとき
$\sup_{x\in A}\lVert f(x)-u(x) \rVert<\eta<r$から
$u(x), f(x)\in A_{r}$なので
$f=\pi\circ f$と組み合わせて
補題\ref{lem:retractnameraka}から
$x\in A$
のとき
\[
\rVert f(x)-h(x)\lVert =\rVert \pi \circ f(x)-\pi \circ v(x)\lVert\le M_{r}\lVert f(x)- v(x)\lVert\le 2M_{r}\eta
\]
である．
よって
$\sup_{x\in A}\lVert f(x)- h(x)\lVert \le 2M_{r}\eta<\epsilon$となり，
証明が終わる．
\end{proof}



\begin{thm}[球面の拡張手性と次元の関係性スペシフィック]\label{thm:specifickakutyou}
$n\in \zz_{\ge 0}$
とし
$X$を
$\rr^{n}$
のコンパクト部分集合とし
$A$を$X$の閉部分集合とする．
このとき任意の連続写像
$f\colon A\to \mysphere{n}$
は拡張となる連続写像
$g\colon X\to \mysphere{n}$
を持つ．
\end{thm}
\begin{proof}
これはちゃんと証明する．
まず
系\ref{cor:extspnameraka}
を使って
(なめらかな)関数
$h\colon \rr^{n}\to \mysphere{n}$
で
あって
\begin{enumerate}
\item 
$\sup_{x\in A}\lVert f(x) -h(x) \rVert<1$. 
\item 
像$h(\rr^{n})$は$0\in \rr^{n+1}$
を含まない．
\end{enumerate}
を満たすものをとる．
次に
ティーチェの拡張定理を使って
$f$の拡張となる連続関数
$w\colon X\to \rr^{n+1}$
を得る．
そして
\[
V=\{\, x\in X\mid \lVert h(x)-w(x) \rVert<1\, \}
\]
と定義すると$V$は$X$の開集合である．
ここで$A$上で
$\lVert h(x)-f(x)\rVert <1$となるという仮定から
$A\subset V$である．
さてウリゾーンの補題から
連続写像$\phi\colon X\to [0, 1]$
で
$A$上で$1$を常にとり
$X\setminus V$上で
恒等的に$0$となるものが存在する．
ここで$J\colon X\times [0, 1]\to \rr^{n+1}$
を
\[
J(x, t)=(1-t)h(x)+tw(x)
\]
と定義するとこれは連続写像になる．
今から
$J(x, \phi(x))$が任意の$x\in X$
についてノンゼロであることを証明しよう．

まず$x\in V$については$h$の値域は
$\mysphere{n}$なので
任意の$t\in [0, 1]$について
以下の計算を得る．
\begin{align*}
\lVert J(x, t)\rVert&\ge 
\lVert h(x)\lVert -\lVert h(x)-h(x, t) \rVert\\
&=\lVert h(x)\lVert -\lVert th(x)-tw(x) \rVert\\
&=\lVert h(x)\lVert -t\lVert h(x)-w(x) \rVert\\
&\ge \lVert h(x)\lVert -1\cdot \lVert h(x)-w(x) \rVert\\
&>\lVert h(x)\lVert -1\\
&=0
\end{align*}
つまり
$J(x, t)\neq 0$
である．
特に$J(x, \phi(x))\neq 0$である．

次に$x\not \in V$の時は$\phi(x)$が$0$の値を取り続けることから
恒等的に$J(x, \phi(x))=h(x)$
である．特に$J(x, \phi(x))\neq 0$
である．
以上から任意の$x\in X$
で$J(x, \phi(x))\neq 0$となることがわかった．
そこで$g\colon X\to \mysphere{n}$
を
\[
g(x)=\frac{J(x, \phi(x))}{\lVert J(x, \phi(x))\rVert}
\]
と定義すれば
$x\in A$上では確かに$f$と一致する．
これで証明が終わる．
\end{proof}

\begin{thm}\label{thm:26thismainthm}
$n\in \zz_{\ge 0}$
とし
$A$を
$\mysphere{n}$
の閉部分集合とする．
このとき任意の連続写像
$f\colon A\to \mysphere{n}$
は拡張となる連続写像
$g\colon \mysphere{n}\to \mysphere{n}$
が存在する．
\end{thm}
\begin{proof}
$\mysphere{n}$
を上半球$\myhemispp{n}$
と
下半球$\myhemispm{n}$
に分ける．
この二つの共通部分である赤道を
$\myeqsp$
とおく．
そして
$S=\myhemispm{n}\cap A$
で
$T=\myhemispp{n}\cap A$
とする．
このとき，
まず
$\myhemispm{n}$
について
$\myhemispm{n}$
は$\mydisc{n}$
に
同相なので
定理
\ref{thm:specifickakutyou}
を適応して
$f|_{S}\colon S\to \mysphere{n}$
を
$g_{0}\colon \myhemispm{n}\to \mysphere{n}$
に拡張する．
このとき$g_{0}$は$S$上で$f$と同じ値を取るので
$\myhemispm{n}$と$A$の共通部分
$S$で貼り合わさって
$f$は$\myhemispm{n}\cup A$上まで拡張されたことになる．この拡張された関数を
$g_{1}\colon \myhemispm{n}\cup A\to \mysphere{n}$
とする．
次に$\myhemispp{n}$について
$g_{1}$の定義域$\cap \myhemispm{n}\cup A$
と$\myhemispp{n}$
の共通部分は
$\myeqsp\cup T$
である．
ここで$\myhemispp{n}$
と
$\myeqsp\cup T$, 
そして
$g_{1}|_{\myeqsp\cup T}\colon \myeqsp\cup T\to \mysphere{n}$
に
定理
\ref{thm:specifickakutyou}
を適応すると
連続写像
$g_{2}\colon \myhemispp{n}\to \mysphere{n}$
をえる．
もちろん
$g_{1}$と$g_{2}$の定義域$\myhemispp{n}$
と$\myhemispm{n}$
の共通部分
$\myeqsp$
上では$g_{1}$と$g_{2}$
は同じ値を取るので
連続写像が貼り合わさって
連続写像$g\colon \mysphere{n}\to \mysphere{n}$
をえる．
作り方から
$g|_{A}=f$なので
定理が示された．
\end{proof}


\section{領域不変性定理}

以下の定理により
ユークリッド空間内のコンパクト部分集合$K$の境界点は
$K$の位相的言明のみで特徴付できることがわかる．
これが領域不変性定理の肝である．

この説の内容は
\cite{dempa3}と同じである．
テレンス・タオ\cite{tao}のブログでは
こうした境界点の位相的特徴付けも簡略化することにより
簡単な領域不変性定理の証明をえているが，
この定理自体をなくしてしまうことは勿体無いと思ったので，この文書では簡略化などはしていない．
\begin{thm}\label{thm:26junbidaiji}
$n\in \zz_{\ge 1}$
とし
$K$を$\rr^{n}$
のコンパクト部分集合とする．
このとき$p\in \partial_{\rr^{n}} K$
となるための必要十分条件は
$p$の$K$内の任意の近傍
$N$について
$p$
の$K$内の開近傍$U$
であって$U\subset N$を満たしかつ
任意の連続写像
$f\colon K\setminus U\to \mysphere{n-1}$
は連続な拡張$g\colon K\to \mysphere{n-1}$
を持つようなものが存在することである．
\end{thm}
\begin{proof}
まず最初に$p\in \partial K$を仮定して定理の条件を導く．
$N$を$p$の任意の近傍とする．
そして$B$を$p$を中心とする十分小さい半径を持つ開球で
$K\cap B\subset N$
を満たすものとする．
このとき$U=K\cap B$
が条件を満たすことを見よう．
ここで$L=\partial B$とすると$L$は
$\mysphere{n-1}$
と同相である．
今ここで任意の連続写像
$f\colon K\setminus U\to \mysphere{n-1}$
を考える．
ここで$K\setminus U$は
$K$の閉集合である．
特に$\rr^{n}$の閉集合でもある．
ここで
$\dim_{T}L\le n-1$
なので
$\dim_{T}(K\cap L)\le n-1$
である．
ここで$L\cap K\subset K\setminus U$
に注意して
定理
\ref{thm:26thismainthm}
から連続写像
$f|_{L\cap K}\colon L\cap K\to \mysphere{n-1}$
の拡張である連続写像
$u\colon L\to \mysphere{n-1}$
を得る．
さて$L$と$K\setminus U$の共通部分は$L\cap K$
で，この上では$u$と$f$は同じ値を取るので写像が貼り合わさって連続写像
$h\colon (K\setminus U)\cup L\to \mysphere{n-1}$
を得る．
そして$p$が$K$
の境界なので
$B$の点であって$K$に属さないもの$q$を取ることができる．
任意の点$x\in K$について
$q$から$x\in K$へ伸ばした直線と
$L$との交点を$b(x)$とする．
このとき$b\colon X\to L$
は連続写像である．
そして$g\colon X\to \mysphere{n-1}$を
\begin{align*}
g(x)
=
\begin{cases}
h(b(x)) & \text{$x\in U\cup (L\cap K)$}\\
f(x) & \text{$x\in K\setminus U$}
\end{cases}
\end{align*}
と定義すると
これら二つの関数$h(b(x))$と$f(x)$は$L\cap K$上で値が一致するのでちゃんと張り合って連続関数になる．
この$g$はちゃんと$f$の拡張になっている．

逆に$p\in K$が
定理の条件を満たすとしよう．
背理法のために$p\in K$は$K$
の内点であるとする．
ここで$p$を中心とする十分小さい半径を持つ
開球$B$で$\cl B\subset K$となるものを考える．
すると定理の仮定より$p$の開近傍$U$で
$U\subset B$
で任意の連続関数$f\colon K\setminus U\to \mysphere{n-1}$が$K$上に拡張するものが存在する．
ここで写像
$f\colon K\setminus U\to \mysphere{n-1}$
を以下のように定義する．
点$p$から$x\in K\setminus U$へと直線を伸ばし
$\partial B$との交点を$f(x)$とする．
これは連続写像で$\partial B$は$\mysphere{n-1}$
と同相なので
ちゃんと連続写像
$f\colon K\setminus U\to \mysphere{n-1}$
を得る．このとき
$f$は$\partial B$上で恒等写像になることに注意しよう．
すると定理の条件より
連続写像$g\colon K\to \mysphere{n-1}$
が存在する．
さて$\cl B\subset K$であった．
このとき$g|_{\cl B}\colon \cl B\to \partial B$
を考えると
$g|_{\cl B}$
は$\cl B$から
$\partial B$
へのレトラクトになる．
これはNo retract theorem に矛盾するので
$p\in K$は境界点でなければならない．
\end{proof}


いよいよ領域不変性定理を証明する．
\begin{thm}[領域不変性定理]\label{thm:maininvdomain}
$n\in \zz_{\ge 1}$
とし
$U$は$\rr^{n}$
の開集合とする．
そして
$f\colon U\to \rr^{n}$
を連続単射とする．
このとき$f(U)$
は$\rr^{n}$の開集合である．
\end{thm}
\begin{proof}
任意の$p\in U$について
$f(p)$が$f(U)$の内点であることを証明すれば良い．
$p$の十分小さい近傍$N$で$\cl N$がコンパクトで
$\cl N \subset U$であるものをとる．
このとき$f|_{\cl N}$を考えると
$f$はコンパクト集合からハウスドルフ空間への
連続単射なので位相的埋め込みになる．
定理\ref{thm:26junbidaiji}より
点$p$が境界点かどうかは$\cl N$の内在的な位相的条件で特徴づけられる．つまり$p$は$\cl N$の内点であるから
$f(p)$は$f(\cl N)$の内点になる．
よって特に$f(p)$は$f(U)$の内点である．
\end{proof}


\begin{cor}
$M$と$N$を同じ次元の境界つき多様体とし
$f\colon M\to N$を同相写像とする．
このとき$f$は$\partial M$を$\partial N$に移す．
\end{cor}
\begin{proof}
もしも$f$が$p\in \partial M$を$N\setminus \partial N$
に移すと，$f^{-1}$を考えることによって局所的には
ユークリッド空間の内点が$p$に写っていることになるが，
$p$は上半平面の境界点なので領域不変性定理に矛盾する．
\end{proof}


%%%%%%%%%%%%%%%%%%%%%%%%%%%%%
%%%%%%%%%%%%%%%%%%%%%%%%%%%%%
%%%%%%%%%%%%%%%%%%%%%%%%%%%%%
%%%%%%%%%%%%%%%%%%%%%%%%%%%%%
%%%%%%%%%%%%%%%%%%%%%%%%%%%%%
%%%%%%%%%%%%%%%%%%%%%%%%%%%%%
%%%%%%%%%%%%%%%%%%%%%%%%%%%%%
%%%%%%%%%%%%%%%%%%%%%%%%%%%%%
%%%%%%%%%%%%%%%%%%%%%%%%%%%%%
%%%%%%%%%%%%%%%%%%%%%%%%%%%%%
%%%%%%%%%%%%%%%%%%%%%%%%%%%%%
%%%%%%%%%%%%%%%%%%%%%%%%%%%%%
%%%%%%%%%%%%%%%%%%%%%%%%%%%%%
%%%%%%%%%%%%%%%%%%%%%%%%%%%%%
%%%%%%%%%%%%%%%%%%%%%%%%%%%%%
%%%%%%%%%%%%%%%%%%%%%%%%%%%%%
%%%%%%%%%%%%%%%%%%%%%%%%%%%%%
%%%%%%%%%%%%%%%%%%%%%%%%%%%%%
%%%%%%%%%%%%%%%%%%%%%%%%%%%%%
%%%%%%%%%%%%%%%%%%%%%%%%%%%%%
%%%%%%%%%%%%%%%%%%%%%%%%%%%%%
%%%%%%%%%%%%%%%%%%%%%%%%%%%%%
%%%%%%%%%%%%%%%%%%%%%%%%%%%%%
%%%%%%%%%%%%%%%%%%%%%%%%%%%%%
%%%%%%%%%%%%%%%%%%%%%%%%%%%%%
%%%%%%%%%%%%%%%%%%%%%%%%%%%%%
%%%%%%%%%%%%%%%%%%%%%%%%%%%%%
%%%%%%%%%%%%%%%%%%%%%%%%%%%%%
%%%%%%%%%%%%%%%%%%%%%%%%%%%%%
%%%%%%%%%%%%%%%%%%%%%%%%%%%%%
%%%%%%%%%%%%%%%%%%%%%%%%%%%%%
%%%%%%%%%%%%%%%%%%%%%%%%%%%%%
%%%%%%%%%%%%%%%%%%%%%%%%%%%%%
%%%%%%%%%%%%%%%%%%%%%%%%%%%%%
%%%%%%%%%%%%%%%%%%%%%%%%%%%%%
%%%%%%%%%%%%%%%%%%%%%%%%%%%%%
%%%%%%%%%%%%%%%%%%%%%%%%%%%%%
%%%%%%%%%%%%%%%%%%%%%%%%%%%%%
%%%%%%%%%%%%%%%%%%%%%%%%%%%%%
%%%%%%%%%%%%%%%%%%%%%%%%%%%%%
%%%%%%%%%%%%%%%%%%%%%%%%%%%%%
%%%%%%%%%%%%%%%%%%%%%%%%%%%%%
%%%%%%%%%%%%%%%%%%%%%%%%%%%%%
%%%%%%%%%%%%%%%%%%%%%%%%%%%%%
%%%%%%%%%%%%%%%%%%%%%%%%%%%%%
%%%%%%%%%%%%%%%%%%%%%%%%%%%%%
%%%%%%%%%%%%%%%%%%%%%%%%%%%%%
%%%%%%%%%%%%%%%%%%%%%%%%%%%%%



\begin{thebibliography}{99}

\bibitem{dempafx}
電波通信
「不動点定理」
https://concious4410.hatenablog.com/entry/2017/08/27/182714



\bibitem{dempa}

電波通信, 「次元論ショートコース：0次元から始める位相次元」, 
https://concious4410.hatenablog.com/entry/2022/03/12/161532


\bibitem{dempa2}
電波通信, 「不動点定理と同値な命題」, 
https://concious4410.hatenablog.com/entry/2026/04/05/181441

\bibitem{dempa3}
電波通信，「次元論からの領域不変性定理の証明」，
https://concious4410.hatenablog.com/entry/2026/04/15/220550

\bibitem{chara}
M. G. Charalambous, 
\textit{Dimension theory. A selection of theorems and counterexamples} Springer/Atlantis Press (2019). 

\bibitem{HW}W. Hurewicz and H. Wallman,\textit{Dimension Theory},Princeton Univ. Press,1948


\bibitem{kihara}
T. Kihara, 
\textit{The Brouwer invariance theorems in reverse mathematics},
Forum of Mathematics, Sigma,
vol.8,  2020, 
 doi: 10.1017/fms.2020.52. 
 
 \bibitem{kurat}
 C. Kuratowski, 
 \textit{Une M\'{e}thode de Prolongement des Ensembles Relativement Ferm\'{e}s ou Ouverts}, Colloq. Math. 1 (1948), 273-278, 
doi:10.4064/cm-1-4-273-278. 
\bibitem{Kul}
W. Kulpa, 
\textit{The Poiancr\'{e}--Miranda Theorem}, 
Amer. Math. Monthly 104, No. 6, 545--550 (1997)
doi: 10.2307/2975081

\bibitem{tao}
T. Tao, 
\textit{Brouwer’s fixed point and invariance of domain theorems, and Hilbert’s fifth problem}, 
https://terrytao.wordpress.com/2011/06/13/brouwers-fixed-point-and-invariance-of-domain-theorems-and-hilberts-fifth-problem/
\bibitem{pears}
 A. R. 
 Pears,
 \textit{Dimension theory of general spaces}, 
 {Cambridge} {University} {Press}
1975. 

\end{thebibliography}




			\end{document}

